\documentclass[10pt,a4paper]{amsart}
\usepackage{amsmath,amssymb,amsthm,graphicx}
% Retained minimal preamble for the Stacks Project chapter more-algebra.tex (More on Algebra).
% The chapter file pinned-original-source.tex carries only \input{preamble} on its first line, so no
% preamble text lives inside the chapter itself.  The Stacks Project uses one single shared preamble for
% every chapter; the snapshot stages it identically as sources/stacks-chow/preamble.tex and
% sources/stacks-etale-cohomology/preamble.tex.  The block below is copied mechanically, line for line,
% from sources/stacks-chow/preamble.tex lines 14-18 (the Xy-pic setup) and 154-217 (the native theorem
% declarations together with the mathematical macros).  Nothing was added, removed or redefined; the
% wrapper class and the amsmath/amssymb/amsthm/graphicx base packages are supplied by the assembler.
% Chapter: more-algebra.tex -- More on Algebra.
% For commutative diagrams we use Xy-pic
\usepackage[all]{xy}

% We use 2cell for 2-commutative diagrams.
\xyoption{2cell}

% Theorem environments.
%
\theoremstyle{plain}
\newtheorem{theorem}[subsection]{Theorem}
\newtheorem{proposition}[subsection]{Proposition}
\newtheorem{lemma}[subsection]{Lemma}

\theoremstyle{definition}
\newtheorem{definition}[subsection]{Definition}
\newtheorem{example}[subsection]{Example}
\newtheorem{exercise}[subsection]{Exercise}
\newtheorem{situation}[subsection]{Situation}

\theoremstyle{remark}
\newtheorem{remark}[subsection]{Remark}
\newtheorem{remarks}[subsection]{Remarks}

\numberwithin{equation}{subsection}

% Macros
%
\def\lim{\mathop{\mathrm{lim}}\nolimits}
\def\colim{\mathop{\mathrm{colim}}\nolimits}
\def\Spec{\mathop{\mathrm{Spec}}}
\def\Hom{\mathop{\mathrm{Hom}}\nolimits}
\def\Ext{\mathop{\mathrm{Ext}}\nolimits}
\def\SheafHom{\mathop{\mathcal{H}\!\mathit{om}}\nolimits}
\def\SheafExt{\mathop{\mathcal{E}\!\mathit{xt}}\nolimits}
\def\Sch{\mathit{Sch}}
\def\Mor{\mathop{\mathrm{Mor}}\nolimits}
\def\Ob{\mathop{\mathrm{Ob}}\nolimits}
\def\Sh{\mathop{\mathit{Sh}}\nolimits}
\def\NL{\mathop{N\!L}\nolimits}
\def\CH{\mathop{\mathrm{CH}}\nolimits}
\def\proetale{{pro\text{-}\acute{e}tale}}
\def\etale{{\acute{e}tale}}
\def\QCoh{\mathit{QCoh}}
\def\Ker{\mathop{\mathrm{Ker}}}
\def\Im{\mathop{\mathrm{Im}}}
\def\Coker{\mathop{\mathrm{Coker}}}
\def\Coim{\mathop{\mathrm{Coim}}}

% Boxtimes
%
\DeclareMathSymbol{\boxtimes}{\mathbin}{AMSa}{"02}

%
% Macros for moduli stacks/spaces
%
\def\QCohstack{\mathcal{QC}\!\mathit{oh}}
\def\Cohstack{\mathcal{C}\!\mathit{oh}}
\def\Spacesstack{\mathcal{S}\!\mathit{paces}}
\def\Quotfunctor{\mathrm{Quot}}
\def\Hilbfunctor{\mathrm{Hilb}}
\def\Curvesstack{\mathcal{C}\!\mathit{urves}}
\def\Polarizedstack{\mathcal{P}\!\mathit{olarized}}
\def\Complexesstack{\mathcal{C}\!\mathit{omplexes}}
% \Pic is the operator that assigns to X its picard group, usage \Pic(X)
% \Picardstack_{X/B} denotes the Picard stack of X over B
% \Picardfunctor_{X/B} denotes the Picard functor of X over B
\def\Pic{\mathop{\mathrm{Pic}}\nolimits}
\def\Picardstack{\mathcal{P}\!\mathit{ic}}
\def\Picardfunctor{\mathrm{Pic}}
\def\Deformationcategory{\mathcal{D}\!\mathit{ef}}

% --- appended 2026-10-12 (Chief): the project own definitions of the reference/slogan/history
% display environments, copied mechanically from the shared Stacks preamble lines 8-12
% (\usepackage{verbatim} plus the three \newenvironment lines, which the chapter preamble keeps
% verbatim from the project preamble). Nothing added, removed or redefined: these environments
% carry no mathematical content and the project itself typesets them as comments. Without these
% three lines a counted statement whose own text uses them fails with Environment slogan undefined.
\usepackage{verbatim}
\newenvironment{reference}{\comment}{\endcomment}
\newenvironment{slogan}{\comment}{\endcomment}
\newenvironment{history}{\comment}{\endcomment}
\usepackage{cleveref}
\crefname{originallemma}{original lemma}{original lemmas}
\Crefname{originallemma}{Original lemma}{Original lemmas}
\makeatletter
\newlabel{more-algebra-lemma-cut-by-koszul-flat}{{0CEQ}{source}{}{sourcepointer-1}{}}
\newlabel{more-algebra-lemma-cut-by-koszul-flat@cref}{{[originallemma][0][]0CEQ}{[1][0][]source}{}{}{}}
\makeatother
\crefname{originallemma}{original lemma}{original lemmas}
\Crefname{originallemma}{Original lemma}{Original lemmas}
\makeatletter
\newlabel{algebra-lemma-flat-localization}{{00HT}{source}{}{sourcepointer-2}{}}
\newlabel{algebra-lemma-flat-localization@cref}{{[originallemma][0][]00HT}{[1][0][]source}{}{}{}}
\makeatother
\crefname{originallemma}{original lemma}{original lemmas}
\Crefname{originallemma}{Original lemma}{Original lemmas}
\makeatletter
\newlabel{algebra-lemma-grothendieck-regular-sequence-general}{{0470}{source}{}{sourcepointer-3}{}}
\newlabel{algebra-lemma-grothendieck-regular-sequence-general@cref}{{[originallemma][0][]0470}{[1][0][]source}{}{}{}}
\makeatother
\crefname{originallemma}{original lemma}{original lemmas}
\Crefname{originallemma}{Original lemma}{Original lemmas}
\makeatletter
\newlabel{lemma-relative-regular-immersion-algebra}{{0CEP}{source}{}{sourcepointer-4}{}}
\newlabel{lemma-relative-regular-immersion-algebra@cref}{{[originallemma][0][]0CEP}{[1][0][]source}{}{}{}}
\makeatother
\crefname{originallemma}{original lemma}{original lemmas}
\Crefname{originallemma}{Original lemma}{Original lemmas}
\makeatletter
\newlabel{lemma-noetherian-finite-all-equivalent}{{09CC}{source}{}{sourcepointer-5}{}}
\newlabel{lemma-noetherian-finite-all-equivalent@cref}{{[originallemma][0][]09CC}{[1][0][]source}{}{}{}}
\makeatother
\title{Flatness after cutting a flat algebra by a quasi-regular sequence}
\author{The Stacks Project Authors}
\date{stacks-project:more-algebra.tex:local-snapshot-20261012; GNU FDL 1.2 or later}
\begin{document}
\maketitle
\setcounter{subsection}{1}
\noindent\textit{Editorial note.} All delivered mathematical text is the author's own, copied from the pinned Stacks Project chapter file more-algebra.tex; no statement or proof step was rewritten, reordered or removed.\par\medskip
\noindent\textit{Editorial note.} The chapter file pinned-original-source.tex carries only \textbackslash{}input\{preamble\} on its first line, so no preamble text lives inside the chapter. The retained file sources/stacks-more-algebra/retained-preamble.tex is a mechanical, line-for-line copy of the project's single shared preamble as staged in the snapshot (Xy-pic setup plus theorem declarations and mathematical macros). Nothing was added, removed or redefined, and the excerpt serves only as preamble, never as a span coordinate.\par\medskip
\noindent\textit{Editorial note.} Both counted spans start at the \textbackslash{}begin line of their environment, so the statement and the author's proof are delivered whole; surrounding prose of the section is deliberately outside the counted spans and is carried as standard prerequisites instead.\par\medskip
\noindent\textit{Editorial note.} The citation fields (source\_id, source\_version, authors, source\_url, license.name, source\_title) are identical to those of every other card this lane has produced from this chapter, so the batch selector treats them as one source.\par\medskip
\noindent\textit{Allowed standard prerequisites (editorial).}\begin{itemize}
\item Flatness is local on the source and the target. Algebra, Lemma \textbackslash{}ref\{algebra-lemma-flat-localization\}, cited at proof lines 8041-8042. Reduces the claim to flatness of a localised ring map.
\item The Grothendieck criterion: a quotient by a sequence is flat when the sequence maps to a regular sequence in the fibre. Algebra, Lemma \textbackslash{}ref\{algebra-lemma-grothendieck-regular-sequence-general\}, cited at proof lines 8043-8045. Converts the required flatness into a regularity statement in the fibre.
\item Quasi-regularity is preserved when passing to the fibre over the residue field. Lemma \textbackslash{}ref\{lemma-relative-regular-immersion-algebra\}, cited at proof lines 8051-8054. Produces the quasi-regular sequence in the fibre of the quotient.
\item A quasi-regular sequence in a Noetherian local ring is a regular sequence, and finite modules cut out by a regular sequence are finite free. Lemma \textbackslash{}ref\{lemma-noetherian-finite-all-equivalent\}, cited at proof lines 8055-8057. Upgrades quasi-regularity in the fibre to regularity, which is what the Grothendieck criterion needs.
\end{itemize}\medskip
\ifdefined\hypertarget\hypertarget{sourcepointer-1}{}\fi
\noindent\textit{Source reference.} Original lemma 0CEQ: focused-seeds/AGP/sources/stacks-snapshot-20261011/tags.txt, source lines 15585--15585. The lemma's own canonical tag, quoted so that readers can resolve the label inside the delivered text; it points at the counted statement itself and carries no proof content.\par\medskip
\ifdefined\hypertarget\hypertarget{sourcepointer-2}{}\fi
\noindent\textit{Source reference.} Original lemma 00HT: focused-seeds/AGP/sources/stacks-snapshot-20261011/tags.txt, source lines 655--655. Navigation pointer to a statement the counted proof refers to but does not contain; nothing load-bearing is delegated to it.\par\medskip
\ifdefined\hypertarget\hypertarget{sourcepointer-3}{}\fi
\noindent\textit{Source reference.} Original lemma 0470: focused-seeds/AGP/sources/stacks-snapshot-20261011/tags.txt, source lines 5491--5491. Navigation pointer to a statement the counted proof refers to but does not contain; nothing load-bearing is delegated to it.\par\medskip
\ifdefined\hypertarget\hypertarget{sourcepointer-4}{}\fi
\noindent\textit{Source reference.} Original lemma 0CEP: focused-seeds/AGP/sources/stacks-snapshot-20261011/tags.txt, source lines 15584--15584. Navigation pointer to a statement the counted proof refers to but does not contain; nothing load-bearing is delegated to it.\par\medskip
\ifdefined\hypertarget\hypertarget{sourcepointer-5}{}\fi
\noindent\textit{Source reference.} Original lemma 09CC: focused-seeds/AGP/sources/stacks-snapshot-20261011/tags.txt, source lines 11825--11825. Navigation pointer to a statement the counted proof refers to but does not contain; nothing load-bearing is delegated to it.\par\medskip
\begin{lemma}
\label{lemma-cut-by-koszul-flat}
Let $A' \to B'$ be a ring map. Let $I \subset A'$ be an ideal.
Set $A = A'/I$ and $B = B'/IB'$. Let $f'_1, \ldots, f'_r \in B'$. Assume
\begin{enumerate}
\item $A' \to B'$ is flat and of finite presentation,
\item $I$ is locally nilpotent,
\item the images $f_1, \ldots, f_r \in B$ form a quasi-regular sequence,
\item $B/(f_1, \ldots, f_r)$ is flat over $A$.
\end{enumerate}
Then $B'/(f'_1, \ldots, f'_r)$ is flat over $A'$.
\end{lemma}

\begin{proof}
Set $C' = B'/(f'_1, \ldots, f'_r)$. We have to show $A' \to C'$ is flat.
Let $\mathfrak r' \subset C'$ be a prime ideal lying over
$\mathfrak p' \subset A'$. We let $\mathfrak q' \subset B'$
be the inverse image of $\mathfrak r'$.
By Algebra, Lemma \ref{algebra-lemma-flat-localization}
it suffices to show that $A'_{\mathfrak p'} \to C'_{\mathfrak q'}$ is flat.
Algebra, Lemma \ref{algebra-lemma-grothendieck-regular-sequence-general}
tells us it suffices to show that $f'_1, \ldots, f'_r$ map to
a regular sequence in
$$
B'_{\mathfrak q'}/\mathfrak p'B'_{\mathfrak q'} =
B_\mathfrak q/\mathfrak p B_\mathfrak q =
(B \otimes_A \kappa(\mathfrak p))_\mathfrak q
$$
with obvious notation. What we know is that $f_1, \ldots, f_r$
is a quasi-regular sequence in $B$ and that $B/(f_1, \ldots, f_r)$
is flat over $A$. By Lemma \ref{lemma-relative-regular-immersion-algebra}
the images $\overline{f}_1, \ldots, \overline{f}_r$
of $f'_1, \ldots, f'_r$ in $B \otimes_A \kappa(\mathfrak p)$ form a
quasi-regular sequence. Since $(B \otimes_A \kappa(\mathfrak p))_\mathfrak q$
is a Noetherian local ring, we conclude by Lemma
\ref{lemma-noetherian-finite-all-equivalent}.
\end{proof}

\end{document}
